\scanpage{102}
[ N.B. \quad \uline{Lemme} \ $X$ préschéma, $U$ ouvert rétrocompact, $f:X'\to X$ morphisme \uline{fid plat} et \uline{qu cpt}, $U'=f^{-1}(X')$, On suppose $\boxed{i'_*(\underline{\mathcal O}_{U'})\xleftarrow{\sim}\underline{\mathcal O}_{X'}}$. (ce qui est \uline{équivalent} à la formule $\boxed{i_*(\deleted{\unsure{U}}\underline{\mathcal O}_U)\xleftarrow{\sim}\underline{\mathcal O}_X}$).

\diagram{\begin{tikzcd}[ampersand replacement=\&,column sep=small,row sep=small]
\underline E' \& {}\\
X'\arrow[d,"f"] \& U'\arrow[l,"i'"']\arrow[d,"g"]\\
X \& U\arrow[l,"i"']\\
{} \& \underline F\unsure{.|{}_{0}}
\end{tikzcd}}

Soit $P_{\unsure{m|n}}(X)$ l'ens des classes, à isom près, de faisceaux loc. libres \deleted{\unsure{des|de}} type fini sur $X$, définissons de $\overline{\mathrm m}$ $P(X'),P(U),P(U')$. Alors le diagramme

\diagram{\begin{tikzcd}[ampersand replacement=\&]
P(X)\arrow[r]\arrow[d] \& P(X')\arrow[d]\\
P(U)\arrow[r] \& P(U')
\end{tikzcd}}
% slanted
\margin{N.B. Inutile de se borner à des faisceaux loc libres. On peut \unsure{énoncer|écrire|encore|prendre}, en prenant les faisceaux $\underline E$ tels que $\operatorname{prof} E_x\geq2$ pour $x\in X-U$, et $\underline E'$ en conséquence, \ldots\ [en loc. noeth]}

\uline{cartésien} \deleted{\unsure{cartésien|cocartésien|cohérent}}. En particulier, le diagramme
\diagram{\begin{tikzcd}[ampersand replacement=\&]
\mathrm{Pic}(X)\arrow[r]\arrow[d] \& \mathrm{Pic}(X')\arrow[d]\\
\mathrm{Pic}(U)\arrow[r] \& \mathrm{Pic}(U')
\end{tikzcd}}
est \uline{cartésien}.

En effet, \deleted{\unsure{si}} le foncteur $\underline E\mapsto i^*(\underline E)$ des $\underline E$ loc libres sur $X$ dans les Mod loc libres sur $U$ est \uline{pleinement fidèle}, donc $\deleted{\unsure{P}}P(X)\to P(U)$ est injectif, [et de $\overline{\mathrm m}$ $P(X')\to P(U')$ l'est] donc $P(X)\to P(U)\times_{P(U')}P(X')$ est injectif. \deleted{Soient d'\unsure{autre|abord}} \inserted{Prouvons la} surjectivité. Soient $\underline F$ loc. libre sur $U$, $\underline E'$ loc libre sur $X'$, tels que $i'{}^*(\underline E')\simeq g^*(\underline F)$. Je dis que $\underline E=i_*(\underline F)$ est loc libre $\underset{\substack{f\ \text{fid plat}\\\text{qu. cpt}}}{\Leftarrow} f^* i_*(\underline F)$ est loc. libre, Or
\[
f^*i_*(\underline F)\underset{f\ \mathrm{plat}}{\simeq} i'_*g^*(\underline F)\underset{g^*(\underline F)\simeq i'{}^*(\underline E')}{\simeq} i'_*i'{}^*(\underline E')\underset{i'_*(\underline{\mathcal O}_{U'})\simeq\underline{\mathcal O}_{X'}}{\simeq}\underline E',\quad\text{OK.}
\]
On a $i^*(\underline E)\simeq\underline F$, Je dis que
\[
 f^*(\underline E)\simeq\underline E'\quad \Longleftarrow\quad
 \underbrace{i'{}^*f^*(\underline E)\simeq i'{}^*(\underline E')}\longrightarrow
 g^*i^*(\underline E)\simeq g^*(\underline F)
 \quad\text{O.K.}
\]
% archive sheet number: 101

